\begin{theorem}[Tashkinov's theorem]
Let $(G,e,\varphi)$ be a $k$-triple and let $T$ be a Tashkinov tree with
respect to $e$ and $\varphi$.  Then $V(T)$ is elementary with respect to
$\varphi$.
\end{theorem}
\begin{proof}
We use the exact notions of a $k$-triple, tree-sequence, Tashkinov tree, and
tree vertex set from \noderef{KTriple}, \noderef{TreeSequence},
\noderef{TashkinovTree}, and \noderef{TreeSequenceVertexSet}; elementariness
means \noderef{ColoredVertexSetElementary}, with missing colors as in
\noderef{MissingColors}.  Write
\[
 T=(y_0,e_1,y_1,\ldots,e_n,y_n),\qquad e_1=e.
\]
The edge identities and vertices displayed here are pairwise distinct by the
tree-sequence data, so no step below identifies parallel edges merely because
their endpoints agree.

We first record the coloring facts used in the argument.  In a proper partial
coloring leaving just $e$ uncolored, the edges of any two fixed colors form
components which are paths, even cycles, or isolated vertices: properness
gives degree at most two, with at most one incident edge of either color at
each vertex.  Swapping the two colors on one component preserves properness.
At an internal vertex of a path the two incident colors are merely exchanged;
at an end, the unique incident one becomes the other color; and no edge
outside the component changes.  We call this a Kempe interchange.  A global
permutation of the color names also preserves properness and every
Tashkinov-witness relation.

The two ends $y_0,y_1$ of $e$ each miss at least two colors: after $e$ is
uncolored, at most $\Delta(G)-1$ colors are present there, while
$k\geq\Delta(G)+1$.  Every other vertex misses at least one color, since at
most $\Delta(G)$ colors are present.  Also $G$ has no proper full
$k$-edge-coloring, because the $k$-triple equality is
$\chi'(G)=k+1$.

Suppose, for a contradiction, that the theorem fails.  Among all proper
partial $k$-edge-colorings of $G$ leaving $e$ uncolored and all non-elementary
Tashkinov trees for those colorings, choose a pair $(\varphi,T)$ of minimum
rank as follows.  For a displayed tree-sequence
$S=(z_0,f_1,z_1,\ldots,f_m,z_m)$, let $\rho(S)$ be the least
$r\in\{1,\ldots,m\}$ for which the suffix
$z_r,f_{r+1},z_{r+1},\ldots,f_m,z_m$ is a path in the displayed order; the
choice $r=m$ always works.  Order pairs lexicographically by
$(\rho(S),m)$.  The natural numbers are well ordered, so a minimum-rank bad
pair exists.  Put $q=\rho(T)$.

Every proper prefix $T_j=(y_0,e_1,y_1,\ldots,e_j,y_j)$ is elementary;
otherwise it would be a bad Tashkinov tree of smaller rank.  It is still a
Tashkinov tree because every edge and every earlier missing-color witness in
the prefix is inherited from $T$.  The same prefix conclusion remains true
after any recoloring that preserves $T$ as a bad Tashkinov tree with the same
rank.

We need two consequences of this minimality.

First, every proper prefix $T_j$ has at least four colors which are missing
on $V(T_j)$ and unused on its colored tree edges $e_2,\ldots,e_j$.  Indeed,
the lower bounds above give at least
\[
  2+2+(j-1)=j+3
\]
missing-color incidences on its $j+1$ vertices.  Prefix elementariness makes
all those colors distinct.  Only $j-1$ colored tree edges occur, so at most
$j-1$ of the missing colors are used on them.  At least four remain unused.

Second, consider $i<j<n$ and distinct colors $\alpha,\beta$ such that
$\alpha$ is missing at $y_i$, $\beta$ is missing at $y_j$, and neither
color is used on $e_2,\ldots,e_j$.  The $(\alpha,\beta)$-component through
$y_j$ is a path ending at $y_i$.  To prove this, observe first that $y_j$
has degree at most one in this two-colored component because it misses
$\beta$.  If the other end were not $y_i$, interchange the two colors on
the component.  The color $\alpha$ would then be missing at both $y_i$ and
$y_j$, while no tree edge or witness in $T_j$ changes.  This would make the
proper prefix $T_j$ non-elementary, a contradiction.  Thus the other end is
$y_i$.

After interchanging on this path, the missing statuses of $\alpha$ and
$\beta$ are exchanged at $y_i,y_j$ and unchanged at every other tree
vertex.  Edges through $T_j$ do not change.  A later tree edge with a color
outside $\{\alpha,\beta\}$ keeps its witness.  If a later witness in those
two colors is lost at one of $y_i,y_j$, the other vertex supplies the
witness; both precede that later edge.  If the later edge itself changes
color, $y_i$ or $y_j$ supplies the corresponding new witness.  Hence $T$
remains Tashkinov.  The multiplicity with which either color is missing on
$V(T)$ is unchanged, so $T$ remains non-elementary and has the same rank.
We call this operation a controlled interchange.  Notice also that a
two-colored path cannot have three distinct vertices at which one of its two
colors is missing: each would have degree at most one, but a path has only
two ends.

We next show $q>2$.  If $n=1$, non-elementariness gives a color missing at
both ends of $e$; assigning that color to $e$ would give a full proper
$k$-edge-coloring, impossible.  Thus $n\geq2$.  If $q\leq2$, exchange the
names of $y_0,y_1$ if necessary; then the whole displayed sequence is a
path.  Since $T_{n-1}$ is elementary but $T$ is not, some color $\beta$ is
missing at $y_n$ and an earlier $y_i$.  If $i<n-1$, choose a color
$\gamma$ missing at $y_{i+1}$.  Both are unused on $T_{i+1}$.  For
$\gamma$, any earlier use would have a still earlier Tashkinov witness and
would contradict prefix elementariness.  For $\beta$, the only possible
witness is $y_i$, and the only new candidate edge is $e_{i+1}$, which cannot
have color $\beta$ because it is incident with the $\beta$-missing vertex
$y_i$.  A controlled interchange moves the missing color $\beta$ from
$y_i$ to $y_{i+1}$ and leaves it missing at $y_n$, the latter because a
two-colored path has only two ends.  Repetition makes $\beta$ missing at
both $y_{n-1}$ and $y_n$.  Recolor $e_n$ with $\beta$, and let $\delta$ be
its former color.  The Tashkinov condition supplies a vertex $y_j$ with
$j\leq n-2$ missing $\delta$, while removal of color $\delta$ from $e_n$
makes it missing also at $y_{n-1}$.  Thus $T_{n-1}$ is a shorter bad
Tashkinov path, a contradiction.  Hence $q>2$.

Because $T_{n-1}$ is elementary, every color repeated as a missing color on
$T$ is missing at $y_n$ and at some earlier vertex.  We now treat three
exhaustive cases.

\textbf{Case 1: $q=n$.}
Then $e_n$ does not join $y_{n-1}$ to $y_n$, for otherwise the path suffix
would start before $y_n$ and $q<n$.  We put the coloring into three useful
forms, always by controlled interchanges, so the minimum-rank conclusions
remain valid.

First, there is a color $\beta$ missing at $y_n$ and some $y_j$ with
$j\leq n-2$, and unused on every tree edge.  Start with any repeated missing
color $\beta$ at $y_j,y_n$.  The first consequence above, applied to
$T_{n-2}$, gives four missing colors unused there.  The two remaining edges
$e_{n-1},e_n$ use at most two of them, so choose a different color
$\alpha$, missing at some $y_m$ with $m\leq n-2$, which is unused on all of
$T$.  If $j=n-1$, prefix elementariness shows that $\beta$ is unused on
$T_{n-1}$, and incidence with the $\beta$-missing $y_n$ shows it is not the
color of $e_n$; a controlled interchange between $y_m$ and $y_{n-1}$ gives
the claim.

Suppose instead that $j\leq n-2$ and that $\beta$ is used.  Follow the
maximal $(\alpha,\beta)$-chain from $y_n$.  If it first meets $T_{n-2}$ at
$y_h$, append the chain segment from $y_h$ to $y_n$ to $T_{n-2}$.  Its
internal vertices occur outside the prefix, and every appended edge has one
of the two colors missing at the retained vertices $y_m,y_j$.  It is thus a
bad Tashkinov tree whose path index is at most $n-1<q$, impossible.  The
chain therefore avoids $T_{n-2}$.  Every tree edge has an end in this prefix,
including $e_n$ because its earlier end is not $y_{n-1}$, so the chain
contains no tree edge.  Interchange on the chain and then globally exchange
the names $\alpha,\beta$.  Tree-edge colors and witnesses are preserved,
and the resulting color called $\beta$ is missing at $y_m,y_n$ and unused
on $T$.  This proves the first normal form.

Second, if $\gamma=\varphi(e_n)$, then $\gamma$ is present at $y_{n-1}$.
Otherwise the Tashkinov witness for $e_n$ and elementariness of $T_{n-1}$
make $y_{n-1}$ the unique vertex of that prefix missing $\gamma$, so
$\gamma$ is unused on $T_{n-1}$.  Apply a controlled
$(\beta,\gamma)$-interchange between the earlier vertex missing the unused
$\beta$ from the first normal form and $y_{n-1}$.  Its path cannot contain
$y_n$, since $y_n$ misses $\beta$ and is incident with the
$\gamma$-colored $e_n$, which would make it a third end.  Thus $e_n$ keeps
color $\gamma$, while $\gamma$ becomes present at $y_{n-1}$.

Third, some color is missing at $y_n$ and at a vertex of $T_{n-2}$.  Choose
a repeated missing color $\beta'$ at $y_{j'},y_n$.  There is nothing to do
if $j'\leq n-2$.  If $j'=n-1$, then $\beta'$ is unused on all of $T$ by
prefix elementariness and its incidence with $y_n$.  Choose a different
color $\alpha'$ missing on $T_{n-2}$ and unused on the whole tree, as in the
first normal form.  A controlled interchange moves $\beta'$ from
$y_{n-1}$ into $T_{n-2}$ without affecting its absence at $y_n$.  Neither
color is $\gamma$, because $\gamma$ is used on $e_n$ and is present at
$y_{n-1}$, so the second normal form remains true.

Delete $y_{n-1}$ and $e_{n-1}$ from the displayed sequence but retain
$e_n,y_n$ as its last pair.  Since $q=n$, the earlier end of $e_n$ lies in
$T_{n-2}$, so this is a tree-sequence.  The second normal form says that the
missing-$\gamma$ witness for $e_n$ is not $y_{n-1}$ and therefore remains
in $T_{n-2}$.  The third normal form makes the shortened tree bad.  Its path
index is at most $n-1<n=q$, contradicting minimality.

\textbf{Case 2: $q<n$, and some color is missing at $y_n$ and at a vertex
$y_i$ with $i\geq q$.}
Let $\beta$ be such a color.  If $i<n-1$, choose $\gamma$ missing at
$y_{i+1}$.  The suffix beginning at $y_q$ is a path, so $e_{i+1}$ joins
$y_i$ and $y_{i+1}$.  The same prefix-elementarity argument used when
proving $q>2$ shows that $\beta,\gamma$ are unused on $T_{i+1}$.  A
controlled interchange advances the second missing occurrence of $\beta$
from $y_i$ to $y_{i+1}$ without changing its absence at $y_n$.  Repetition
makes $\beta$ missing at $y_{n-1},y_n$.  Recolor $e_n$ with $\beta$ and
call its former color $\delta$.  Its earlier missing-color witness lies in
$T_{n-2}$ because $\delta$ was present at the incident vertex $y_{n-1}$.
After the recoloring, $\delta$ is missing also at $y_{n-1}$.  Hence
$T_{n-1}$ is a bad Tashkinov tree of smaller rank, a contradiction.

\textbf{Case 3: $q<n$, and no color missing at $y_n$ is missing at a
vertex $y_i$ with $i\geq q$.}
There is a repeated missing color $\beta$ at $y_n$ and a vertex $y_i$ with
$i<q$.  We may arrange $i\leq q-2$.  If initially $i=q-1$, apply the
four-unused-colors conclusion to $T_{q-2}$.  Even after allowing for
$e_{q-1}$, choose a color $\alpha\ne\beta$ missing at some
$y_m$, $m\leq q-2$, and unused on $T_{q-1}$.  Prefix elementariness also
makes $\beta$ unused there.  A controlled interchange moves the earlier
missing occurrence of $\beta$ to $y_m$.

Choose $\gamma$ missing at $y_q$.  It is unused on $T_q$, since any use on
a tree edge would have an earlier missing witness and violate prefix
elementariness.  Moreover $\gamma\ne\beta$ by the defining assumption of
this case.  Applying the four-unused-colors conclusion to $T_{q-2}$ and
allowing for $e_{q-1},e_q$, choose a color $\varepsilon\ne\gamma$, missing
at some $y_{m'}$ with $m'\leq q-2$, which is unused on $T_q$.

If $\varepsilon$ is missing at $y_n$, a controlled
$(\varepsilon,\gamma)$-interchange between $y_{m'}$ and $y_q$ moves that
missing occurrence to $y_q$ while leaving it at $y_n$.  This is Case 2,
already impossible.  Thus suppose $\varepsilon$ is present at $y_n$, so in
particular $\varepsilon\ne\beta$.  Follow the maximal
$(\varepsilon,\beta)$-chain from $y_n$.  If it first meets $T_{q-1}$ at
$y_h$, append the chain segment from $y_h$ to $y_n$ to $T_{q-2}$ when
$h\leq q-2$, and to $T_{q-1}$ when $h=q-1$.  In either case both alternating
colors have missing witnesses in the retained prefix.  The resulting tree
is bad because $\beta$ remains missing at the retained $y_i$ and at $y_n$,
and its path index is at most $q-1$, a contradiction.

The chain therefore avoids $T_{q-1}$.  Interchange
$\varepsilon,\beta$ on it.  No edge of $T_q$ changes, because every such
edge has an end in $T_{q-1}$.  Witnesses through $T_q$ are unchanged, and a
later affected edge still has the earlier missing-$\beta$ witness at $y_i$
or the earlier missing-$\varepsilon$ witness at $y_{m'}$.  The tree remains
Tashkinov, and now $\varepsilon$ is missing at both $y_{m'}$ and $y_n$.
The colors $\varepsilon,\gamma$ remain unused on $T_q$, with $\gamma$
missing at $y_q$.  A controlled interchange between $y_{m'}$ and $y_q$
therefore moves the repeated missing color to $y_q,y_n$, again reducing to
the impossible Case 2.

The three cases exhaust $q=n$ and $q<n$, and in the latter situation they
exhaust whether a missing color of $y_n$ repeats at or after $y_q$.  Every
case contradicts the minimum choice.  No bad pair exists, so the original
$V(T)$ is elementary with respect to $\varphi$.
\end{proof}
